% !TeX program = xelatex
% Companion to the Lean sources in AIV5/mmrank21; see README.md in this directory.
\documentclass[11pt,a4paper]{article}
\usepackage[margin=25mm,headheight=15pt]{geometry}
\usepackage{fontspec}
\setmainfont{DejaVu Serif}[Scale=0.93]
\setsansfont{DejaVu Sans}[Scale=0.93]
\setmonofont{DejaVu Sans Mono}[Scale=0.8]
\usepackage[russian]{babel}
\usepackage{amsmath,amssymb,amsthm,mathtools}
\usepackage{tikz-cd}
\usepackage{microtype}
\usepackage{enumitem}
\usepackage{needspace}
\usepackage{booktabs,longtable}
\usepackage[unicode,colorlinks=true,linkcolor=blue!45!black,urlcolor=blue!45!black]{hyperref}
\hypersetup{pdftitle={Нижняя оценка 21 для умножения матриц 3 x 3},
pdfauthor={Andrey Ivanov},
pdfsubject={Русский математический конспект доказательства}}
\usepackage{fancyhdr}
\pagestyle{fancy}\fancyhf{}\fancyhead[L]{\small Нижняя оценка $R_K\langle3,3,3\rangle\geqslant21$}\fancyhead[R]{\thepage}
\setlength{\parindent}{1.3em}\setlength{\parskip}{3pt}
\setlist{itemsep=3pt,topsep=5pt}
\emergencystretch=2em
% Формулы внутри абзацев не разрываются после операций и отношений.
\binoppenalty=10000
\relpenalty=10000
\interdisplaylinepenalty=10000
\numberwithin{equation}{section}
\theoremstyle{plain}
\newtheorem{theorem}{Теорема}[section]
\newtheorem*{maintheorem}{Основная теорема}
\newtheorem{lemma}[theorem]{Лемма}
\newtheorem{prop}[theorem]{Предложение}
\newtheorem{cor}[theorem]{Следствие}
\theoremstyle{definition}
\newtheorem{definition}[theorem]{Определение}
\newtheorem{remark}[theorem]{Замечание}
\renewcommand{\proofname}{Доказательство}
\DeclareMathOperator{\tr}{tr}\DeclareMathOperator{\rk}{rk}
\DeclareMathOperator{\im}{im}\DeclareMathOperator{\supp}{supp}
\DeclareMathOperator{\Mat}{Mat}\DeclareMathOperator{\Span}{span}
\DeclareMathOperator{\Hom}{Hom}\DeclareMathOperator{\diag}{diag}
\newcommand{\kk}{K}\newcommand{\one}{\mathbf1}
\newcommand{\pair}[2]{\langle #1,#2\rangle}
\newcommand{\col}{\mathcal C}\newcommand{\test}{\mathcal T}
\newcommand{\rankone}[2]{#1#2}
\newcommand{\A}{\mathcal A}\newcommand{\B}{\mathcal B}\newcommand{\N}{\mathcal N}
\newcommand{\E}{\mathcal E}\newcommand{\Sig}{\Sigma}
\newcommand{\leqs}{\leqslant}\newcommand{\geqs}{\geqslant}
\begin{document}
\shorthandoff{"}
\begin{center}
{\LARGE Нижняя оценка для умножения матриц}\\[6pt]
{\Large $R_K\langle3,3,3\rangle\geqslant21$}\\[8pt]
{\normalsize Математический конспект}
\end{center}
\begin{quote}\small
Это сгенерированный ИИ человекочитаемый конспект Lean-доказательства
из \href{https://github.com/AIV5/mmrank21}{AIV5/mmrank21}, а не научная
статья. Текст не проходил независимой математической проверки или
рецензирования. Новизна и приоритет результатов и доказательств
не заявляются. Формальным источником утверждений и доказательств
служит Lean-код; его проверка не удостоверяет корректность этого
изложения. Сопроводительная карта связывает утверждения с Lean-модулями.
\end{quote}
\begin{maintheorem}
Пусть $K$ --- произвольное поле. Если для всех матриц $A,B\in\Mat_3(K)$
выполняется тождество
\[
 BA=\sum_{i=1}^{r} a_i(A)b_i(B)W_i,
\]
где $a_i,b_i$ --- линейные формы, а $W_i\in\Mat_3(K)$ фиксированы,
то $r\geqs21$.
\end{maintheorem}

\tableofcontents
\newpage
\section{Матричная запись и основные тождества}

Все векторы пространств $P,Q,S$ считаются строками; матрицы действуют
справа. Временно допускаем размеры $p=\dim P$, $q=\dim Q$, $s=\dim S$,
где $p,s\geqs1$, $q\geqs2$.
Направления коэффициентов и входов таковы:
\[
\begin{tikzcd}[column sep=large]
 P\arrow[r,"U_i"] & Q\arrow[r,"V_i"] & S\arrow[ll,bend left=35,"W_i"]
\end{tikzcd}
\qquad
\begin{tikzcd}[column sep=large]
 S\arrow[r,"B"]&Q\arrow[r,"A"]&P\arrow[ll,bend left=35,"C"]
\end{tikzcd}
\]
Например, $U_i$ имеет формат $p\times q$, $A$ --- $q\times p$,
$BA$ --- $s\times p$. Произведение $UV$ означает сначала действие $U$,
затем действие $V$. Столбец $\eta$ задаёт форму $x\mapsto x\eta$ на строках.
Для столбца $a$ и строки $b$ матрица $ab$ имеет ранг один, если оба
множителя ненулевые. Это обычное матричное произведение.
Прямые $Ka$ и $Kb$ будем называть соответственно её столбцовым
и строковым направлениями.
\emph{Столбцовый носитель} семейства матриц --- линейная оболочка всех
их столбцов; \emph{строковый носитель} --- оболочка всех строк,
то есть сумма образов при правом действии.

\begin{definition}
\emph{Билинейный алгоритм длины $r$} состоит из $r$ троек матриц
$(U_i,V_i,W_i)$, индексированных множеством $I$, и удовлетворяет тождеству
\begin{equation}\label{eq:algorithm}
 BA=\sum_{i\in I}\tr(U_iA)\tr(V_iB)W_i\qquad\text{для всех }A,B.
\end{equation}
Алгоритмы здесь точные: тождество выполняется над самим полем,
без предельного перехода. Далее это уточнение опускается.
Минимальную длину алгоритма указанного формата обозначаем
$R_K\langle p,q,s\rangle$.
\end{definition}
Любая линейная форма на матрицах единственным образом имеет вид
$A\mapsto\tr(UA)$: значение формы на матричной единице $E_{ab}$ равно
$U_{ba}$. Поэтому это определение эквивалентно формулировке теоремы.

\begin{definition}
Матрицу $C:P\to S$, которой сопоставляется выходная форма
$BA\mapsto\tr(CBA)$, будем называть \emph{тестом}.
\emph{Отображения, заданные коэффициентными формами}, обозначаются
\[
 \alpha(A)_i=\tr(U_iA),\quad
 \beta(B)_i=\tr(V_iB),\quad
 \gamma(C)_i=\tr(W_iC).
\]
Их значения лежат в $K^I$. В этом пространстве используем покоординатное
произведение $xy=(x_i y_i)_i$, единицу $\one=(1)_i$ и спаривание
$\pair{x}{y}=\sum_i x_i y_i$. Для подпространств $X,Y\subseteq K^I$
через $XY$ обозначаем \emph{линейную оболочку} всех произведений $xy$.
Носитель вектора $x$ есть $\supp x=\{i:x_i\ne0\}$; носитель
подпространства --- объединение носителей его векторов.
\end{definition}
Для невырожденного спаривания \emph{аннулятором} подпространства
называем множество всех элементов противоположного пространства,
спаривающихся с ним в нуль. Для матриц используется спаривание
$\tr(XY)$, для координатных векторов --- $\pair{x}{y}$.
В последнем случае пишем также $X^\perp$ и говорим об ортогональности;
никакой положительной определённости это не предполагает.

\begin{lemma}\label{lem:identities}
Отображения $\alpha,\beta,\gamma$ инъективны. Для всех допустимых матриц
выполняются тождества
\begin{align}
 \tr(CBA)&=\sum_i\alpha(A)_i\beta(B)_i\gamma(C)_i,\label{eq:trace}\\
 CB&=\sum_i\beta(B)_i\gamma(C)_iU_i,&
 AC&=\sum_i\gamma(C)_i\alpha(A)_iV_i,\label{eq:cycles}\\
 \sum_i\gamma(C)_i U_i X V_i&=\tr(X)C
       &&(X\in\Mat_q(K)),\label{eq:sandwich}\\
 \sum_i\tr(XV_iTU_i)W_i&=\tr(X)T
       &&(T\in\Mat_{s,p}(K)).\label{eq:dual-sandwich}
\end{align}
\end{lemma}
\begin{proof}
Первое равенство получается умножением \eqref{eq:algorithm} на $C$ и
взятием следа. Невырожденность спаривания $\tr(XY)$, проверяемая на
матричных единицах, даёт оба равенства \eqref{eq:cycles}.
Если, например, $\alpha(A)=0$, то $BA=0$ для всех $B$ и потому $A=0$.
Циклическая перестановка даёт остальные инъективности.

Для \eqref{eq:sandwich} сначала подставим в \eqref{eq:trace} матричные
единицы вместо $A,B$. Получаем равенство коэффициентов
\[
 \sum_i(U_i)_{a b}(V_i)_{c d}\gamma(C)_i
       =\delta_{bc}C_{ad}.
\]
Умножение на $X_{bc}$ и суммирование даёт \eqref{eq:sandwich}.
Следовая двойственность по $C$ даёт \eqref{eq:dual-sandwich}.
Все равенства полиномиальны в коэффициентах и сохраняются при расширении поля.
\end{proof}

Нулевые слагаемые можно удалять. Предполагая существование алгоритма
длины не более двадцати, перейдём к алгебраическому замыканию $K$ и удалим
нулевые слагаемые. Если не оговорено иное, дальше, до заключительного
раздела, $K$ алгебраически замкнуто, все $U_i,V_i,W_i$ ненулевые
и $r\leqs20$.

\section{Координатные произведения и кондукторы}

\begin{definition}
Для подпространств $M,N\subseteq K^I$ их \emph{кондуктор} есть
\[
 (N:M)=\{d\in K^I:md\in N\text{ для каждого }m\in M\}.
\]
Это подпространство $K^I$; двоеточие не обозначает факторпространство.
\emph{Мультипликативный стабилизатор} $X$ есть алгебра
$(X:X)=\{a:aX\subseteq X\}$. Её элементы не обязаны быть обратимыми.
\end{definition}

\begin{lemma}\label{lem:nz-growth}
Пусть $X,Y\subseteq K^I$ ненулевые и $xy=0$, $x\in X$, $y\in Y$
влечёт $x=0$ или $y=0$. Тогда
\[
 \dim XY\geqs\dim X+\dim Y-1.
\]
\end{lemma}
\begin{proof}
Индукция по $\dim X$. При $\dim X=1$ умножение на ненулевой $x\in X$
инъективно на $Y$. Иначе выберем координату $i$, на которой какое-нибудь
произведение ненулевое. Пространства
$X_0=\{x\in X:x_i=0\}$ и $(XY)_0=\{v\in XY:v_i=0\}$ имеют
коразмерность один в $X$ и $XY$. При этом $X_0Y\subseteq(XY)_0$.
Индукционное предположение для $X_0,Y$ даёт требуемую оценку.
\end{proof}

\begin{definition}
Для $G\subseteq I$ положим
\[
 X[G]=\{x\in X:x_i=0\text{ при }i\notin G\},\qquad
 X|_G=\{(x_i)_{i\in G}:x\in X\}.
\]
Это соответственно \emph{координатное сечение} и \emph{ограничение}
пространства. Они, вообще говоря, имеют разные размерности.
\end{definition}

\begin{lemma}\label{lem:two-conductors}
Пусть $M=\Span(\one,x,y)$, $N,B\subseteq K^I$ и $MB\perp N$.
Тогда
\[
 \dim(N:M)+\dim(B:M)\geqs3\dim N+3\dim B-2r.
\]
\end{lemma}
\begin{proof}
Для любого $L\subseteq K^I$ рассмотрим
\[
 \Phi_L:L^3\longrightarrow K^I\oplus K^I,\qquad
 (l_0,l_1,l_2)\longmapsto(l_1-xl_0,l_2-yl_0).
\]
Его ядро состоит из $(l_0,xl_0,yl_0)$ с $l_0\in(L:M)$.
Следовательно,
\[
 \dim\im\Phi_L=3\dim L-\dim(L:M).
\]
На $K^I\oplus K^I$ форма
\[
 \omega((a,b),(c,d))=\pair{a}{d}-\pair{b}{c}
\]
невырождена и альтернирована, то есть $\omega(v,v)=0$, в том числе
в характеристике два. Образы $\Phi_N$ и $\Phi_B$ ортогональны:
после раскрытия скобок члены с множителями $1,x,y$ равны нулю по условию,
а члены с $xy$ сокращаются. Сумма размерностей ортогональных
подпространств не превосходит $2r$. Подставим найденные размерности.
\end{proof}

\subsection{Связность и координатные компоненты}

\begin{definition}
Конечное непустое семейство ненулевых векторов $(v_i)_{i\in G}$ называется
\emph{связным}, если при каждом разбиении $G=G_1\sqcup G_2$ на две
непустые части
\[
 \Span\{v_i:i\in G_1\}\cap\Span\{v_i:i\in G_2\}\ne0.
\]
Разбивая несвязные семейства, получаем единственное самое мелкое
разбиение, для которого оболочка семейства является прямой суммой
оболочек частей. Части называются компонентами связности.
\end{definition}
Единственность следует также из описания через минимальные линейные
зависимости: индексы одной такой зависимости должны лежать в одной
части; транзитивное замыкание этих связей и даёт компоненты.

\begin{lemma}\label{lem:connected}
Пусть ненулевые матрицы ранга один $a_i b_i$ образуют связное семейство.
Тогда
\[
 \dim\Span\{a_i\}+\dim\Span\{b_i\}
       \leqs\dim\Span\{a_i b_i\}+1.
\]
\end{lemma}
\begin{proof}
Начнём с одной матрицы. Пусть для уже выбранной части $J$ неравенство
доказано. Если $J\ne G$, связность позволяет выбрать минимальную
непустую часть $H\subseteq G\setminus J$, оболочка которой пересекает
оболочку $J$. Матрицы части $H$ независимы; пересечение одномерно и
порождено $T=\sum_{i\in H}c_i a_i b_i\ne0$, где все $c_i\ne0$.
Иначе можно убрать один индекс из $H$.

Обозначим прежние оболочки факторов через $A_0,B_0$, а увеличенные ---
через $A_1,B_1$. Столбцы $T$ лежат в $A_0$, а строки --- в $B_0$.
Отсюда следует
\[
 (\dim A_1-\dim A_0)+(\dim B_1-\dim B_0)\leqs |H|-1.
\]
Вот непосредственно линейно-алгебраическая проверка последнего шага.
Матрица $T$ факторизуется как $XY$, где столбцы $X$ равны $c_i a_i$,
а строки $Y$ равны $b_i$. Проецируя $X$ по $A_0$, получаем
$\bar X Y=0$; проецируя строки по $B_0$, получаем $X\bar Y=0$.
Если сумма двух приростов равнялась бы $|H|$, то
$\im\bar Y=\ker\bar X$; из $X\bar Y=0$ следовало бы, что $X$ зануляет
$\ker\bar X$. Поэтому $XY=0$, поскольку $\bar X Y=0$.
Это противоречит $T\ne0$. Обычная оценка ранга произведения даёт
сначала сумму приростов не больше $|H|$, так что остаётся $|H|-1$.
Оболочка матриц при присоединении $H$ увеличивается ровно на $|H|-1$.
Неравенство сохраняется; конечная индукция завершает доказательство.
\end{proof}

\begin{definition}
\emph{Координатные компоненты} подпространства $X\subseteq K^I$ ---
компоненты связности ненулевых координатных форм $x\mapsto x_i$ на $X$.
Нулевые координатные формы не включаем. Эквивалентно, это минимальные
непустые блоки $G$, для которых
$X=\bigoplus_G X[G]$.
Такие компоненты называем \emph{активными}: на каждой из их координат
пространство $X$ не тождественно нулевое.
\end{definition}

\begin{lemma}\label{lem:component-cost}
Пусть $\one\in M$, $D=(N:M)$, $E=MD$. Для каждой активной координатной
компоненты $G$ пространства $E$ имеем
\[
 \dim(M|_G)+\dim D[G]\leqs\dim E[G]+1.                 \tag{*}
\]
Если $\dim E=\dim D+1$, найдётся такая компонента с $\dim(M|_G)=2$.
Если $E=D$, то $D$ порождено общими собственными векторами всех
операторов $w\mapsto mw$, $m\in M$.
\end{lemma}
\begin{proof}
Координатные проекторы, сохраняющие $E$, сохраняют и $D$:
$m\pi_G(d)=\pi_G(md)\in E\subseteq N$. Поэтому
$E[G]=(M|_G)D[G]$. Выберем базисы в $M|_G$ и $D[G]$.
Координаты произведения задаются билинейными формами
$(m,d)\mapsto m_i d_i$; их матрицы имеют ранг один.
Линейные связи между этими матрицами в точности совпадают со связями
между координатными формами на $E[G]$, поскольку произведения
порождают $E[G]$. Поэтому семейство матриц связно, а размерность его
оболочки равна $\dim E[G]$. Первые и вторые факторы порождают
пространства размерностей $\dim(M|_G)$ и $\dim D[G]$: это полные
наборы координатных форм на соответствующих пространствах.
Лемма~\ref{lem:connected} даёт (*).

Сумма приростов по компонентам равна общему приросту. Из (*) каждый
$\dim(M|_G)$ не больше двух при приросте один. Если все они равны одному,
всякое $m$ постоянно на каждом блоке (в $M$ есть единица), поэтому
$ME=E$. Отсюда $E\subseteq(N:M)=D$, противоречие.

При $MD=D$ имеем $M\subseteq(D:D)$: все множители принадлежат
мультипликативному стабилизатору $D$.
Ограничения диагональных операторов $w\mapsto mw$ на $D$
одновременно диагонализуемы. Это можно увидеть без теоремы о спектре:
разобьём $I$ по равенству всех значений $m_i$; проектор на каждый класс
выражается полиномом от конечного числа этих диагональных операторов
с помощью интерполяции. Он сохраняет $D$. На каждом полученном сечении
все операторы скалярны. Любой его ненулевой вектор --- общий собственный.
\end{proof}

\section{Нормировка алгоритма и размерностные оценки}

\subsection{Выбор гиперплоскости}
Для матричного подпространства $H\subseteq\Mat_{p,q}(K)$ положим
\[
 \col(H)=\{a\in K^{p\times1}:ab\in H
                 \text{ для всех }b\in K^{1\times q}\}.
\]
Это подпространство столбцов, а не ядро матриц из $H$.

\begin{lemma}\label{lem:generic}
Пусть над бесконечным полем задан алгоритм формата $\langle p,q,s\rangle$
произвольной длины с ненулевыми $U_i,V_i,W_i$, где $p,s\geqs1$, $q\geqs2$.
Тогда можно выбрать ненулевой столбец $\eta\in K^{q\times1}$, гиперплоскость
$F=\{x:x\eta=0\}\subset Q$ и ненулевые масштабы коэффициентов так,
чтобы для
\[
 M=\{\alpha(\eta u):u\in P\},\quad
 B_F=\{\beta(B):\im B\subseteq F\},\quad N=\im\gamma
\]
выполнялись следующие свойства:
\begin{enumerate}
\item $\one\in M$; пространства $M,B_F,N$ имеют полный носитель;
\item $a\in M$, $b\in\im\beta$, $ab=0$ влекут $a=0$ или $b=0$;
\item для каждого $G\subseteq I$, $i\in G$, если
$U_i\eta\in\col(\Span\{U_j:j\in G\})$, то все столбцы $U_i$
принадлежат этому же подпространству.
\end{enumerate}
Здесь полный носитель пространства означает, что ни одна координата
не равна нулю тождественно на всём пространстве.
\end{lemma}
\begin{proof}
Условие в пункте 3 получается исключением конечного числа собственных
линейных подпространств: для фиксированных $G,i$ рассмотрим прообраз
$\col(\Span U_G)$ при $\eta\mapsto U_i\eta$. Если прообраз не весь,
выберем $\eta$ вне него.

Кроме того, исключим $\ker U_i$ и те $\eta$, для которых $F\subseteq\ker V_i$.
Последнее возможно лишь при $\rk V_i=1$ и исключает одну прямую
направлений $\eta$. Это обеспечивает полный носитель $M,B_F$.
Полный носитель $N$ следует из $W_i\ne0$.

Для каждого $J$ с $\Span V_J\ne\Mat_{q,s}(K)$ выберем ненулевую
матрицу $B_J$, следово ортогональную всем $V_i$, $i\in J$,
и дополнительно потребуем $B_J\eta\ne0$.
Пусть $u\ne0$, $J=\supp\alpha(\eta u)$. Если $\Span V_J$ неполна,
тождество алгоритма даёт $B_J\eta u=0$, противоречие.
Значит, $V_J$ порождают всё пространство. Если теперь
$\alpha(\eta u)\beta(B)=0$, то $\beta(B)$ равна нулю на $J$,
откуда $B=0$. Получен пункт 2.

Над бесконечным полем конечное объединение собственных линейных
подпространств не покрывает всё пространство: произведение их ненулевых
линейных уравнений --- ненулевой полином. Так существует нужная $\eta$.
Затем выбираем $u_0$ с $c_i=\tr(U_i\eta u_0)\ne0$ для всех $i$ и заменяем
$U_i$ на $c_i^{-1}U_i$, $W_i$ на $c_iW_i$. Алгоритм и ранги сохраняются,
а $\alpha(\eta u_0)=\one$. Все перечисленные свойства сохраняются при
этих масштабированиях.
\end{proof}

Алгоритм вместе с выбором $\eta,u_0$, удовлетворяющим этой лемме,
далее называем \emph{нормированным}.
Пишем $m(u)=\alpha(\eta u)$ и определяем
\[
 D=(N:M),\qquad E=MD,\qquad Z=(B_F:M),\qquad
 d=\dim D,\quad e=\dim E,\quad z=\dim Z.
\]
Число $g=e-d$ называется \emph{приростом размерности}.
Оно неотрицательно, поскольку $\one\in M$ даёт $D\subseteq E$.
Для $p=q=s=3$ имеем
\[
 \dim M=3,\qquad\dim B_F=6,\qquad\dim N=9,\qquad MB_F\perp N.
\]
Последняя ортогональность следует из $B\eta=0$ и \eqref{eq:trace}.
Лемма~\ref{lem:two-conductors} сразу даёт
\begin{equation}\label{eq:dz}
 d+z\geqs45-2r.
\end{equation}

\subsection{Матричные представления кондукторов}

\begin{definition}
Для $d_0\in D$, $z_0\in Z$, $u\in P$ определим матрицы $C_{d_0}(u)$
и $B_{z_0}(u)$ однозначными равенствами
\[
 \gamma(C_{d_0}(u))=m(u)d_0,\qquad
 \beta(B_{z_0}(u))=m(u)z_0,\qquad \im B_{z_0}(u)\subseteq F.
\]
Существование следует из определения кондуктора, единственность ---
из инъективности коэффициентных отображений. Зависимость от каждого
аргумента линейна.
\end{definition}
\[
\begin{tikzcd}[column sep=huge,row sep=large]
 P\arrow[r,"u\mapsto C_{d_0}(u)"]\arrow[d,"m"']&
 \Mat_{p,s}(K)\arrow[d,"\gamma"]\\
 M\arrow[r,"a\mapsto ad_0"']&N
\end{tikzcd}
\qquad
\begin{tikzcd}[column sep=huge,row sep=large]
 P\arrow[r,"u\mapsto B_{z_0}(u)"]\arrow[d,"m"']&
 \{B:\im B\subseteq F\}\arrow[d,"\beta"]\\
 M\arrow[r,"a\mapsto az_0"']&B_F
\end{tikzcd}
\]

\begin{lemma}\label{lem:reciprocity}
Для всех $u,v\in P$ выполнены
\begin{align}
 vC_{d_0}(u)&=uC_{d_0}(v),\label{eq:reciprocity}\\
 C_{d_0}(u)B_{z_0}(v)&=C_{d_0}(v)B_{z_0}(u).\label{eq:compatibility}
\end{align}
Пусть $\E=\gamma^{-1}(E)$. Для всякого $C\in\E$
\begin{equation}\label{eq:symmetrizer}
 uCB_{z_0}(v)=vCB_{z_0}(u).
\end{equation}
Кроме того, $B_{z_0}(u)=0$ влечёт $z_0=0$ или $u=0$.
\end{lemma}
\begin{proof}
В \eqref{eq:trace} подставим $A=\eta v$ и
$\gamma(C)=m(u)d_0$. Правая часть симметрична по $u,v$;
варьируя $B\eta$, получаем \eqref{eq:reciprocity}.
Аналогично произведение $m(u)m(v)d_0z_0$ и следовая двойственность
дают \eqref{eq:compatibility}. Для $C=C_{d_0}(w)$ получаем
\[
 uC_{d_0}(w)B_{z_0}(v)
 =wC_{d_0}(u)B_{z_0}(v)
 =wC_{d_0}(v)B_{z_0}(u)
 =vC_{d_0}(w)B_{z_0}(u).
\]
Такие $C$ порождают $\E$, поэтому верно \eqref{eq:symmetrizer}.
Последнее утверждение следует из пункта 2 леммы~\ref{lem:generic}.
\end{proof}

\begin{lemma}\label{lem:physical-bound}
Для нормированного алгоритма формата $\langle p,q,s\rangle$
произвольной длины $r$ справедливо
\[
 ps+(q-1)e\leqs r.
\]
В частности, при $p=q=s=3$, $r\leqs20$ имеем $e\leqs5$.
\end{lemma}
\begin{proof}
Рассмотрим $r$ линейных отображений
$Q_i:T\mapsto V_iTU_i\eta$ из $\Mat_{s,p}(K)$ в пространство столбцов $K^q$.
Из \eqref{eq:sandwich} и \eqref{eq:reciprocity} следуют два равенства
\begin{align*}
 \sum_i\gamma(C)_iQ_i(T)&=\tr(TC)\eta,\\
 \sum_i(d_0)_i\alpha(\theta u)_iQ_i(T)
       &=\tr(TC_{d_0}(u))\theta.
\end{align*}
Для второго достаточно проверить $T=av$ ранга один: его левая часть
после \eqref{eq:trace} равна
$(vC_{d_0}(u)a)\theta$, что и требуется; такие матрицы порождают всё
пространство. Таким образом, оболочка $Q_i$ содержит все отображения
$T\mapsto\tr(TC)\eta$ и все $T\mapsto\tr(TC)\theta$ с $C\in\E$,
$\theta\in K^q$. Эти два пространства имеют размерности $ps$ и $qe$,
а их пересечение --- размерность $e$. Значит, их сумма имеет размерность
$ps+(q-1)e\leqs r$.
\end{proof}

\Needspace{10\baselineskip}
\begin{cor}\label{cor:two-two-n}
Для любого поля $K$ и целого $n\geqs1$ выполняется
\[
 R_K\langle2,2,n\rangle\geqs\left\lceil\frac{7n}{2}\right\rceil.
\]
\end{cor}
\begin{proof}
Алгоритм над $K$ сохраняется при расширении поля до алгебраического
замыкания, а удаление нулевых слагаемых не увеличивает длину.
Поэтому достаточно оценить длину $r$ алгоритма без нулевых
слагаемых над бесконечным полем. Ограничение $r\leqs20$
здесь не предполагается.

Применим лемму~\ref{lem:generic} при $p=q=2$, $s=n$.
Инъективность коэффициентных отображений даёт
\[
 \dim M=2,\qquad \dim N=2n,\qquad \dim B_F=n.
\]
Поскольку $\one\in M$, выберем $x$ так, что $M=\Span(\one,x)$.
Из $MB_F\perp N$ и коммутативности покоординатного произведения
следует $MN\perp B_F$. Невырожденность спаривания даёт
\[
 \dim MN\leqs r-n.
\]
Рассмотрим сюръективное линейное отображение
\[
 \Phi:N\oplus N\longrightarrow MN=N+xN,\qquad
 (u,v)\longmapsto xu-v.
\]
Его ядро состоит из пар $(u,xu)$, где $u\in N$ и $xu\in N$.
Это в точности условие $u\in D=(N:M)$. По формуле размерностей
ядра и образа получаем
\[
 d=\dim D=4n-\dim MN\geqs5n-r.
\]
Так как $\one\in M$, имеем $D\subseteq MD$, то есть $e\geqs d$.
Теперь лемма~\ref{lem:physical-bound} даёт
\[
 r\geqs2n+e\geqs2n+d\geqs7n-r.
\]
Следовательно, $2r\geqs7n$; целочисленность $r$ завершает доказательство.
\end{proof}

\subsection{Совместные симметризаторы}

\begin{definition}
Пусть $B_z(u):S\to F$ зависит билинейно от $z\in Z$, $u\in P$.
\emph{Пространством совместных симметризаторов} назовём
\[
 \Sig(B)=\{C:P\to S:\ uCB_z(v)=vCB_z(u)
                          \text{ для всех }u,v,z\}.
\]
В выбранных базисах это пространство решений линейных условий:
определённые произведения матриц должны быть симметричными.
По лемме~\ref{lem:reciprocity} $\E\subseteq\Sig(B)$.
\end{definition}

\begin{lemma}\label{lem:rank-one-space}
Если все матрицы некоторого линейного пространства имеют ранг не больше
одного, то у них либо общий одномерный столбцовый носитель, либо общий
одномерный строковый носитель.
\end{lemma}
\begin{proof}
Выберем ненулевую $ab$. Для другой ненулевой матрицы $cd$ сумма
$ab+cd$ имеет ранг два, если одновременно независимы $a,c$ и $b,d$.
Поэтому каждый элемент имеет первый фактор, пропорциональный $a$,
или второй, пропорциональный $b$. Если встречаются оба типа с
непропорциональными свободными факторами, их сумма снова имеет ранг два.
Следовательно, один из факторов общий для всего пространства.
\end{proof}

\begin{lemma}\label{lem:sym-small}
Пусть $\dim P=\dim S=3$, $\dim F=2$, а $B_z(u)=0$ влечёт
$z=0$ или $u=0$. Если $\dim Z\geqs3$ и
$\dim\Span\{B_z(u):z,u\}\geqs5$, то $\dim\Sig(B)\leqs1$.
\end{lemma}
\begin{proof}
Пусть $C\in\Sig(B)$ вырождена и $0\ne u\in\ker C$. Из симметрии
следует, что все $B_z(u)$ зануляют $\im C$. Отображение
$z\mapsto B_z(u)$ инъективно, поэтому
\begin{equation}\label{eq:sym-capacity}
 \dim Z\leqs(3-\rk C)\dim F.
\end{equation}
Значит, вырожденный симметризатор имеет ранг не больше одного.

Если есть обратимый симметризатор $C_0$ и ненулевой симметризатор
ранга один $C$, выберем строку $v$ с $vC\ne0$ и
$x=(vC)C_0^{-1}$. Для $u\in\ker C$ два закона симметрии дают
\[
 uC_0B_z(x)=xC_0B_z(u)=vCB_z(u)=uCB_z(v)=0.
\]
Поэтому все $B_z(x)$ зануляют $(\ker C)C_0$, пространство размерности два.
Инъективность по $z$ даёт $\dim Z\leqs2$, противоречие.
Таким образом, при наличии $C_0$ всякий ненулевой элемент $\Sig(B)$
обратим. Два независимых таких элемента невозможны: у $C_0^{-1}C_1$
есть собственное значение, и $C_1-\lambda C_0$ --- ненулевой
вырожденный элемент.

Остаётся случай, когда все симметризаторы имеют ранг не больше одного.
Предположим, что их пространство $H$ имеет размерность хотя бы два.
По лемме~\ref{lem:rank-one-space} возможны два случая.
Если все $C\in H$ имеют общий образ в прямой $L\subset S$, то
напишем два независимых элемента как $a_1s,a_2s$, где $Ks=L$.
Симметрия означает, что отображение $u\mapsto sB_z(u)$ пропорционально
каждой из двух независимых форм $u\mapsto ua_1$, $u\mapsto ua_2$.
Следовательно, оно нулевое: $LB_z(u)=0$ для всех $z,u$.
Поэтому оболочка всех $B_z(u)$
имеет размерность не больше $\dim\Hom(S/L,F)=4$.
Если все $C\in H$ имеют общее ядро-гиперплоскость $H_0\subset P$,
то для $x\notin H_0$ отображение $C\mapsto xC$ инъективно.
Его образ имеет размерность $\dim H\geqs2$. Для ненулевого
$u\in H_0$ симметрия заставляет все $B_z(u)$ занулять этот образ;
снова $\dim Z\leqs(3-\dim H)2\leqs2$. Оба случая невозможны.
\end{proof}

\begin{lemma}\label{lem:five-sym}
В условиях предыдущей леммы, кроме ограничений на $\dim Z$ и оболочку
$B$, предположим $Z\ne0$ и $\dim\Sig(B)\geqs5$.
Тогда существуют обратимая матрица $L$ и инъективное линейное
отображение $f:Z\to F$, такие что
\begin{equation}\label{eq:factor}
 B_z(u)= (Lu^{\mathsf t}) f(z).
\end{equation}
Иными словами, значение на строке $s$ равно $(sLu^{\mathsf t})f(z)$.
\end{lemma}
\begin{proof}
Сначала установим вспомогательный факт. Если независимые матрицы
$A,B\in\Mat_3(K)$ удовлетворяют условию
$Ak=Bk=0\Rightarrow k=0$ для столбцов $k$, то пространство матриц $T$,
для которых $TA$ и $TB$ симметричны, имеет размерность не больше четырёх.
Двойственные к этим условиям уравнения задаются отображением
\[
 (k,l)\longmapsto A[k]_\times+B[l]_\times,\qquad
 [k]_\times=\begin{pmatrix}0&-k_3&k_2\\k_3&0&-k_1\\-k_2&k_1&0\end{pmatrix}.
\]
Здесь $k,l$ --- столбцы; используем обозначение
$k\times l=[k]_\times l$ для указанной формулы над полем $K$.
Если значение равно нулю, умножение справа на $k$ и $l$ показывает,
что $A(k\times l)=B(k\times l)=0$. Значит, $k,l$ зависимы.
Ядро отображения, рассматриваемое как пространство матриц $(k\ l)$
формата $3\times2$, целиком имеет ранг не больше одного.
Если его размерность хотя бы два, лемма~\ref{lem:rank-one-space}
даёт общий первый или второй фактор. В первом случае ядро содержит
$(k,0),(0,k)$ с $k\ne0$; тогда $A[k]_\times=B[k]_\times=0$,
что противоречит нулевому общему ядру. Во втором случае оно имеет вид
$(ak,bk)$ с фиксированными $(a,b)\ne(0,0)$; ненулевая матрица
$aA+bB$ может занулять $[k]_\times$ лишь для $k$ на одной прямой.
Это видно по любой её ненулевой строке: она должна быть пропорциональна $k$.
Поэтому ядро имеет размерность не больше одного, ранг --- не меньше пяти.
Искомое пространство симметризаторов имеет коразмерность хотя бы пять.
Здесь используются нулевые диагонали $[k]_\times$, так что деления
на два нет.

Выберем $z\ne0$ и запишем две выходные координаты билинейного
отображения $(u,s)\mapsto sB_z(u)$ матрицами $A,B$ размера $3\times3$.
Запись имеет вид $(sAu^{\mathsf t},sBu^{\mathsf t})$.
Поэтому $Ak=Bk=0$ влечёт $k=0$: отображение $u\mapsto B_z(u)$ инъективно.
Пять симметризаторов и доказанный факт вынуждают $A,B$ быть
пропорциональными. Их общий ненулевой множитель $L$ обратим.
Сравним его с каждой координатной матрицей для остальных $z$:
общее ядро опять нулевое, и тот же факт даёт пропорциональность $L$.
Коэффициенты пропорциональности составляют линейное отображение $f$.
Условие отсутствия нулевых пар даёт инъективность $f$.
\end{proof}

\begin{lemma}\label{lem:injective-product}
Если семейство второго кондуктора имеет вид \eqref{eq:factor}, то
$\dim NZ=9z$ и $r\geqs9z+3$.
То же заключение верно, если существуют базис строк $(u_j)$ в $P$,
базис строк $(s_j)$ в $S$ и инъективные $f_j:Z\to F$, для которых
\begin{equation}\label{eq:diagonal-factor}
 B_z(u_j)=s_j^*f_j(z).
\end{equation}
Здесь $s_j^*$ --- столбец двойственного базиса: $s_k s_j^*=\delta_{kj}$.
\end{lemma}
\begin{proof}
Пусть $z_1,\ldots,z_z$ --- базис $Z$ и
$\sum_a\gamma(C_a)z_a=0$. Умножим на $m(u)$ и используем определение
$B_{z_a}(u)$, затем \eqref{eq:cycles}. Получим
\[
 \sum_a C_a B_{z_a}(u)=0\qquad\text{для каждого }u.
\]
В случае \eqref{eq:factor} это
$\sum_a C_a Lu^{\mathsf t}f(z_a)=0$.
Независимость строк $f(z_a)$ даёт $C_aLu^{\mathsf t}=0$ для всех $a,u$,
откуда $C_a=0$. В случае \eqref{eq:diagonal-factor} аналогично
$C_a s_j^*=0$ для всех $a,j$, что опять даёт $C_a=0$.
Итак, отображение $(C_a)_a\mapsto\sum_a\gamma(C_a)z_a$ инъективно.
Его образ --- $NZ$, размерность которого равна $9z$.
Так как $Z\subseteq B_F$ и $MB_F\perp N$, имеем $M\perp NZ$.
Следовательно, $r\geqs\dim M+\dim NZ=3+9z$.
\end{proof}

\begin{prop}\label{prop:critical}
Для нормированного алгоритма формата $3\times3\times3$ длины $r\leqs20$
выполняются
\[
 3\leqs d\leqs e\leqs5,\qquad z\leqs2,\qquad d+z\geqs45-2r,
 \qquad e=5\Longrightarrow d\geqs4, z\leqs1.
\]
В частности, $r\geqs20$, а при $r=20$ прирост $g=e-d$ равен нулю или одному.
\end{prop}
\begin{proof}
Из леммы~\ref{lem:nz-growth}, применённой к $M,Z$, следует
$\dim MZ\geqs z+2$, если $Z\ne0$. Так как $MZ\subseteq B_F$,
получаем $z\leqs4$. Вместе с \eqref{eq:dz} это даёт $d\geqs1$.
По лемме~\ref{lem:physical-bound} $e\leqs5$.

Если $z\geqs3$, оболочка матриц $B_z(u)$ имеет размерность
$\dim MZ\geqs5$. Лемма~\ref{lem:sym-small} даёт $e\leqs1$,
поэтому $D=E$ одномерно. Его ненулевой вектор $w=\gamma(C)$ есть общий
собственный вектор умножений на $M$. Выберем активную координату $i$;
пространство $M_0=\{m\in M:m_i=0\}$ имеет размерность два.
Для $m(u)\in M_0$ имеем $C_w(u)=0$.
Из \eqref{eq:compatibility}, выбрав $u_0$ с $m(u_0)=\one$,
получаем $CB_z(u)=0$. Поэтому матричные представители $M_0Z$
лежат в $\Hom(S/\im C,F)$, размерности не больше четырёх.
Лемма~\ref{lem:nz-growth} даёт $z+1\leqs4$, то есть $z\leqs3$.
Теперь $d+z\leqs4$, вопреки \eqref{eq:dz}. Значит, $z\leqs2$, $d\geqs3$.

При $e=5$, $z=2$ леммы~\ref{lem:five-sym} и~\ref{lem:injective-product}
дают $r\geqs21$, противоречие. Поэтому $z\leqs1$, $d\geqs4$.
Если бы $r\leqs19$, то \eqref{eq:dz} дало бы $d+z\geqs7$;
при $e\leqs4$ это невозможно, а при $e=5$ уже $d+z\leqs6$.
Следовательно, $r=20$. Все оставшиеся пары $(d,e)$ имеют разность
ноль или один.
\end{proof}

\section{Носители тестов и атомарные тесты}

\subsection{Две оценки для тестов с заданным носителем}

\begin{definition}
Для $G\subseteq I$ положим
\[
 \test_G=\{C\in\Mat_{p,s}(K):\supp\gamma(C)\subseteq G\},\quad
 \delta_G=\dim\test_G,\quad
 H_G=\sum_{C\in\test_G}\im C,\quad h_G=\dim H_G.
\]
Таким образом, $\test_G$ состоит из тестов с носителем в $G$,
а $H_G$ --- их общий строковый носитель.
\end{definition}

\begin{lemma}\label{lem:budgets}
Пусть столбцы всех $U_i$, $i\in G$, лежат в фиксированном
$k$-мерном пространстве $H\subset K^{p\times1}$, $k<p$, и $m=|G|$.
Тогда столбцы всех $C\in\test_G$ также лежат в $H$, и
\begin{align}
 \delta_G&\leqs k h_G,\label{eq:capacity}\\
 \delta_G+(q-1)h_G&\leqs m,\label{eq:internal}\\
 m+qs+ks&\leqs r+\delta_G.\label{eq:external}
\end{align}
Для \eqref{eq:internal} нужна нормировка $\alpha(\eta u_0)=\one$.
Кроме того,
\begin{equation}\label{eq:mass}
 \Span\{V_i:i\in G\}\supseteq\{V: \im V\subseteq H_G\},
 \qquad \dim\Span V_G\geqs qh_G.
\end{equation}
\end{lemma}
\begin{proof}
Из \eqref{eq:cycles} для $C\in\test_G$ получаем
$CB\in\Span U_G$ и $AC\in\Span V_G$ для всех $A,B$.
Первое утверждение помещает столбцы $C$ в $H$, второе даёт
\eqref{eq:mass}. Формула \eqref{eq:capacity} --- размерность пространства
матриц со столбцами в $H$ и строками в $H_G$.

Пусть $f:K^G\to\Mat_{q,s}(K)$, $f(c)=\sum_{i\in G}c_iV_i$.
Инъективное отображение $\gamma:\test_G\to K^G$ удовлетворяет
$f\gamma(C)=\eta u_0C$. Поэтому
$\dim f(\gamma\test_G)\leqs h_G$. Из неравенства
\[
 \dim\im f+\dim\gamma\test_G\leqs m+\dim f(\gamma\test_G)
\]
и \eqref{eq:mass} следует \eqref{eq:internal}.

Для оценки \eqref{eq:external} ограничим первый вход условием $AH=0$.
Исчезнут все слагаемые $G$. На пространстве тестов со столбцами в $H$
ограничение $\gamma$ на $I\setminus G$ имеет ядро $\test_G$, а значит
ранг $ks-\delta_G$. Выберем столько независимых координатных форм.
Дополняя произвольный тест слагаемым со столбцами в $H$, можно занулить
эти формы; произведение в \eqref{eq:trace} не изменится, поскольку $AH=0$.
Получаем алгоритм для прямоугольного умножения с внутренним размером
$p-k>0$ и с не более чем $r-m-ks+\delta_G$ слагаемыми.
Его коэффициентное отображение на пространстве матриц формата $q\times s$
инъективно, поэтому длина не меньше $qs$. Это и есть \eqref{eq:external}.
\end{proof}

\begin{lemma}\label{lem:rect15}
Билинейное умножение формата $\langle3,2,3\rangle$ требует не меньше
пятнадцати слагаемых.
\end{lemma}
\begin{proof}
Применим нормировку и кондукторы с $p=s=3$, $q=2$.
Тогда $\dim M=3$, $\dim B_F=3$, $\dim N=9$.
Если $Z\ne0$, лемма~\ref{lem:nz-growth} даёт $z\leqs1$;
если $Z=0$, эта оценка тем более верна. Следовательно,
$d\geqs36-2r-z\geqs35-2r$.
Лемма~\ref{lem:physical-bound} даёт $9+e\leqs r$, а $d\leqs e$.
Значит, $44\leqs3r$, откуда $r\geqs15$.
Переход к бесконечному расширению поля и удаление нулевых слагаемых
не увеличивают длину.
\end{proof}

\begin{lemma}\label{lem:point-capacity}
В алгоритме формата $\langle3,3,3\rangle$ длины двадцать не более трёх
матриц $U_i$ имеют столбцы в одной фиксированной прямой. Если таких
индексов ровно три, то $\delta_G=1$, а их $U_i$ порождают всё
трёхмерное пространство матриц со столбцами в этой прямой.
То же верно для строк и для циклических перестановок $U,V,W$.
\end{lemma}
\begin{proof}
Здесь $k=1$. У тестов с носителем в $G$ фиксирован первый фактор, поэтому
$\delta_G=h_G$ и \eqref{eq:mass} даёт $3\delta_G\leqs m$.
В доказательстве \eqref{eq:external} после удаления $G$ и
$3-\delta_G$ дополнительных
членов остаётся формат $\langle3,2,3\rangle$; по лемме~\ref{lem:rect15}
\[
 15\leqs20-m-3+\delta_G,\qquad m\leqs\delta_G+2.
\]
Вместе эти оценки дают $\delta_G\leqs1$, $m\leqs3$.
При $m=3$ получаем $\delta_G=1$.
Для ненулевого $C=as\in\test_G$ матрицы $CB$, где $B$ произвольна,
порождают всё пространство $\{ab:b\in K^{1\times3}\}$ и по
\eqref{eq:cycles} лежат в оболочке $U_G$. Это доказывает утверждение
об оболочке.
Транспонирование и циклическая перестановка \eqref{eq:trace} доказывают
остальные варианты.
\end{proof}

\begin{lemma}\label{lem:partial}
Пусть из произведения матриц $2\times n$ и $n\times2$ вычисляется
линейное пространство выходных форм размерности $\delta\geqs3$.
Тогда число слагаемых не меньше $3n+\delta-3$.
\end{lemma}
\begin{proof}
Сначала $\delta=3$. Трёхмерное пространство тестовых матриц $2\times2$
является следовым аннулятором одной ненулевой матрицы. Независимыми
заменами двух базисов её ранг приводится к одному из двух случаев.
При ранге один вычисляются три формы
$x u^{\mathsf t},x v^{\mathsf t},y u^{\mathsf t}$, где $x,y,u,v\in K^n$.
Формы первого входа на $(0,y)$ отделяют все $y$: иначе третья выходная
форма тождественно занулялась бы. Выберем $n$ независимых из них и
линейной подстановкой $y=Lx$ занулим их. Останутся не более $r-n$
членов, вычисляющих первые два выхода. Они отделяют все $2n$ координат
$(u,v)$, поэтому $r-n\geqs2n$.

При ранге два выходы имеют вид
$x u^{\mathsf t}-y v^{\mathsf t},x v^{\mathsf t},y u^{\mathsf t}$.
На $K^n\oplus K^n$ их матрицы равны
\[
 S_0=\begin{pmatrix}I&0\\0&-I\end{pmatrix},\quad
 S_1=\begin{pmatrix}0&I\\0&0\end{pmatrix},\quad
 S_2=\begin{pmatrix}0&0\\I&0\end{pmatrix}.
\]
У блочного отображения
\[
 (v_0,v_1,v_2)\longmapsto
 (S_1v_2-S_2v_1,\ S_2v_0-S_0v_2,\ S_0v_1-S_1v_0)
\]
ядро нулевое: из последних двух уравнений выразим $v_1,v_2$ через
$v_0$, а первое даёт
$(S_1S_0^{-1}S_2-S_2S_0^{-1}S_1)v_0=0$; оператор в скобках равен
$-I$. Значит, ранг равен $6n$. Одно билинейное слагаемое имеет
во всех трёх $S_a$ одну и ту же матрицу ранга один с тремя скалярными
коэффициентами. В блочном отображении его образ поэтому имеет размерность
не больше ранга соответствующей альтернированной матрицы $3\times3$,
то есть не больше двух. Следовательно, $2r\geqs6n$.
Знак не мешает характеристике два.

При $\delta>3$ коэффициентные формы на пространстве тестов инъективно
вкладывают его в $K^r$. Выберем $\delta-3$ независимых координатных
форм и возьмём их общее ядро, имеющее размерность три. На нём исчезают
$\delta-3$ слагаемых, а только что доказанная оценка даёт
$r-(\delta-3)\geqs3n$.
\end{proof}

\begin{lemma}\label{lem:plane}
В алгоритме длины двадцать не более семи $U_i$ имеют столбцы в одной
плоскости $H\subset K^3$. Для любого множества таких индексов
$\delta_G\leqs2$.
\end{lemma}
\begin{proof}
Сначала нормируем алгоритм по лемме~\ref{lem:generic}; ненулевые
масштабирования не меняют носителей и рассматриваемых размерностей.
Из леммы~\ref{lem:budgets} при $k=2$ получаем
$m\leqs\delta_G+5$, $\delta_G+2h_G\leqs m$,
$\delta_G\leqs2h_G$. Поэтому $h_G\leqs2$.
Если $\delta_G\geqs3$, то $h_G=2$. Тесты из $\test_G$ ограничивают
произведение $2\times3$ на $3\times2$ и вычисляют $\delta_G$ независимых
выходных форм. Лемма~\ref{lem:partial} даёт $m\geqs\delta_G+6$,
противоречие. Поэтому $\delta_G\leqs2$ и $m\leqs7$.
\end{proof}

\subsection{Атомарные тесты}

\begin{definition}
\emph{Атомарный тест} --- ненулевая матрица ранга один $C=as$, где
$a$ --- столбец, $s$ --- строка, такая что
$|\supp\gamma(C)|=3$. Множество $G=\supp\gamma(C)$ называется его
носителем; $Ka$ и $Ks$ --- его столбцовым и строковым направлениями.
\end{definition}

\begin{lemma}\label{lem:atom}
Для любого ненулевого теста $C$ имеем
$|\supp\gamma(C)|\geqs3\rk C$.
Для атомарного теста $C=as$ на его носителе $G$ существуют строки
$b_i$ и столбцы $c_i$ такие, что
\[
 U_i=ab_i,\qquad V_i=c_i s,
 \qquad \gamma(C)_i\, b_jc_i=\delta_{ji}\quad(i,j\in G).
\]
В частности, все $U_i,V_i$ на $G$ имеют ранг один.
Носители двух атомарных тестов либо совпадают, либо не пересекаются.
В первом случае тесты пропорциональны. Столбцовые направления любых трёх
тестов с различными носителями независимы, как и строковые направления.
\end{lemma}
\begin{proof}
Билинейная форма $(A,B)\mapsto\tr(CBA)$ имеет матричный ранг
$3\rk C$: после приведения $C$ к ранговой нормальной форме это сумма
$3\rk C$ попарно независимых координатных произведений.
Представление \eqref{eq:trace} содержит лишь $|\supp\gamma(C)|$
форм ранга один. Это доказывает оценку.

При $C=as$ и ровно трёх членах достигается равенство рангов.
Факторизация билинейной формы полного ранга через $K^3$ показывает,
что первые коэффициентные формы образуют базис форм $A\mapsto bAa$,
а вторые --- базис форм $B\mapsto sBc$. Поэтому $U_i=ab_i$,
$V_i=c_i s$. Равенство формы матричному скалярному произведению даёт
указанную взвешенную двойственность.

Если два атома имеют общий индекс, его ненулевая $U_i$ фиксирует
общее столбцовое направление. Лемма~\ref{lem:point-capacity} допускает
в этом классе лишь три индекса, так что носители совпадают.
Пространство тестов с этим носителем одномерно, значит тесты
пропорциональны. Наконец, зависимость столбцовых направлений трёх различных
атомов поместила бы девять их $U_i$ в одну плоскость, вопреки
лемме~\ref{lem:plane}. Для строковых направлений применим симметрию
с транспонированием.
\end{proof}

\begin{lemma}\label{lem:stationary}
Пусть $0\ne d_0\in D$, $G=\supp d_0$ и $a=\dim(Md_0)$.
Тогда существует $a$-мерное пространство столбцов $H$, для которого
\[
 \Span U_G=\{T:\text{все столбцы }T\text{ лежат в }H\}.
\]
В частности, $\dim\Span U_G=3a$.
Если $w\in N\setminus\{0\}$ --- общий собственный вектор умножений
на $M$, то $w=\gamma(C)$ для атомарного теста $C$.
\end{lemma}
\begin{proof}
Пусть $H$ --- оболочка столбцов всех $C_{d_0}(u)$.
Из взаимности \eqref{eq:reciprocity} следует
$H=(\ker(u\mapsto C_{d_0}(u)))^\perp$: строка $v$ аннулирует все эти
столбцы тогда и только тогда, когда $uC_{d_0}(v)=0$ для всех $u$.
Поэтому $\dim H=\dim Md_0=a$.

Все $C_{d_0}(u)$ имеют носитель в $G$; из \eqref{eq:cycles} следует
$H\subseteq\col(\Span U_G)$. Равенство
$\gamma(C_{d_0}(u))_i=(uU_i\eta)(d_0)_i$ показывает, что
$U_i\eta\in H$ для $i\in G$: обе стороны зануляются на ядре семейства $C_{d_0}$.
Пункт 3 леммы~\ref{lem:generic} теперь помещает все столбцы $U_i$ в
$\col(\Span U_G)$. Это последнее пространство само лежит в $H$,
поскольку его элементы можно получить вычислением матриц $\Span U_G$
на $\eta\ne0$. Следовательно, все столбцы $U_i$ лежат в $H$,
а противоположное включение полного матричного пространства уже
содержится в определении $\col$.

Для общего собственного вектора $w$ имеем $w\in D$, $\dim Mw=1$.
Значит, все $U_i$ на его носителе имеют одно столбцовое направление.
По лемме~\ref{lem:point-capacity} носитель содержит не больше трёх
индексов, а по лемме~\ref{lem:atom} --- не меньше трёх.
Его тест имеет ранг один и является атомарным.
\end{proof}

\section{Исключение положительного прироста}

\subsection{Локальная геометрия растущей компоненты}

\begin{lemma}\label{lem:growing-block}
Если $e=d+1$, существует координатная компонента $G$ пространства $E$
со следующими свойствами:
\begin{enumerate}
\item $\dim(M|_G)=2$, $\dim D[G]=1$,
$E[G]=N[G]$ и $\dim E[G]=2$;
\item $6\leqs |G|\leqs7$; оболочка $U_G$ --- всё пространство матриц
со столбцами в некоторой плоскости;
\item пространство $\test_G$ двумерно, его столбцовый и строковый
носители двумерны; в $\test_G$ есть матрица ранга два;
\item в $G$ не более двух матриц $U_i$ с любым фиксированным столбцовым
направлением и не более двух --- с любым фиксированным строковым
направлением;
\item $N[G]$ не содержит общего собственного вектора умножений на $M$.
\end{enumerate}
\end{lemma}
\begin{proof}
Лемма~\ref{lem:component-cost} даёт компоненту с ограничением размерности два.
В $D[G]$ есть вектор $d_0$ с носителем ровно $G$: носители $D[G]$ и
$E[G]$ совпадают, и над бесконечным полем можно избежать всех
координатных гиперплоскостей. По лемме~\ref{lem:stationary}
оболочка $U_G$ --- всё пространство матриц со столбцами в некоторой
плоскости. Его размерность равна шести.
Значит, $|G|\geqs6$. Лемма~\ref{lem:plane} даёт $|G|\leqs7$ и
$\dim N[G]\leqs2$. Теперь
\[
 2+\dim D[G]\leqs\dim E[G]+1,\qquad
 1\leqs\dim D[G],\qquad E[G]\subseteq N[G]
\]
вынуждают размерности $1,2,2$. Если в $N[G]$ есть общий собственный
вектор $w$, он принадлежит $D[G]$ и порождает его; тогда $MD[G]=Kw$,
вопреки $\dim E[G]=2$.

Столбцовый носитель $\test_G$ двумерен по доказательству
леммы~\ref{lem:stationary}. Из оценки \eqref{eq:internal}
$2+2h_G\leqs7$ получаем $h_G\leqs2$.
Покажем, что $h_G\ne1$. Если все тесты имеют строки в прямой $L$,
то $\test_G=\{as:a\in H\}$ с $\dim H=2$, $Ks=L$.
Рассмотрим отображения
\[
 f_i:\test_G\to\Mat_{3,3}(K),\qquad f_i(C)=\gamma(C)_iV_i.
\]
Их оболочка содержит все $C\mapsto AC$; последние образуют пространство
размерности $3\dim H=6$. После проекции значений по пространству
матриц со строками в $L$ эти шесть измерений исчезают. Поэтому оболочка
отображений $C\mapsto\gamma(C)_i[V_i]$ имеет размерность не больше
$|G|-6\leqs1$.
Для индексов с $[V_i]\ne0$ все ненулевые формы
$\gamma(\cdot)_i|_{\test_G}$ тогда пропорциональны: это следует
вычислением двух пропорциональных отображений на любом тесте, где
одно из них ненулевое. Найдётся ненулевой $C\in\test_G$, зануляющий
все эти формы. Его носитель состоит только из индексов со строками
$V_i$ в $L$; таких индексов не больше трёх по
лемме~\ref{lem:point-capacity}. Поэтому $C$ атомарен.
Для атома $U_i=ab_i$; нормировка даёт
$(u_0a)b_i\eta=1$ на его носителе, откуда
$m(u)_i=(ua)/(u_0a)$ постоянно на нём. Значит, $\gamma(C)$ ---
общий собственный вектор, что уже исключено.

Теперь оба носителя двумерны. Если все тесты имели бы ранг не больше
одного, лемма~\ref{lem:rank-one-space} сделала бы один носитель
одномерным. Поэтому есть тест ранга два.

Три индекса с общим столбцовым направлением имеют ненулевой тест
с носителем в этой тройке по
лемме~\ref{lem:point-capacity}. Это атом и общий собственный вектор,
противоречие. Для строкового направления применим ту же лемму после
транспонирования: три матрицы с общим строковым направлением
должны порождать трёхмерное пространство всех матриц с этим
направлением. Но при дополнительном ограничении столбцов плоскостью
такое пространство имеет размерность только два. Получено противоречие.
\end{proof}

\subsection{Замена одной билинейной пары}

\begin{definition}
Первая сумма рангов алгоритма есть $\sigma_U=\sum_i\rk U_i$.
Из всех алгоритмов длины двадцать над выбранным полем выберем один
с наименьшей $\sigma_U$. Минимум существует, поскольку это непустое
множество натуральных чисел. Нормировка ненулевыми масштабами его сохраняет.
Билинейной парой $U\boxtimes V$ обозначим форму
$(A,B)\mapsto\tr(UA)\tr(VB)$; это не произведение матриц $UV$.
\end{definition}

\begin{lemma}\label{lem:exchange}
В алгоритме с минимальной $\sigma_U$ равенство
\[
 \sum_i a_i(U_i\boxtimes V_i)=X\boxtimes Y,\qquad a_j\ne0,
\]
влечёт $\rk U_j\leqs\rk X$.
\end{lemma}
\begin{proof}
Выразим $U_j\boxtimes V_j$ из равенства и подставим в алгоритм.
Его $j$-й член заменится на пару $(X,Y)$ с выходом $a_j^{-1}W_j$;
для $i\ne j$ выход станет $W_i-a_i a_j^{-1}W_j$.
Все остальные первые коэффициенты сохранятся, длина останется двадцать.
Если $\rk X<\rk U_j$, уменьшится $\sigma_U$.
\end{proof}

\begin{prop}\label{prop:no-growth}
В алгоритме с минимальной первой суммой рангов нормировка
леммы~\ref{lem:generic} удовлетворяет $MD=D$.
\end{prop}
\begin{proof}
По предложению~\ref{prop:critical} достаточно исключить $e=d+1$.
Возьмём компоненту $G$ из леммы~\ref{lem:growing-block}.
Сожмём её два внешних носителя до $K^2$. Первые коэффициенты имеют вид
$U_i=JX_i$, где $J$ --- матрица $3\times2$ ранга два,
$X_i$ --- матрицы $2\times3$. В тестовом пространстве выберем базис
$C,D$ так, чтобы $C$ на двух носителях была единичной матрицей,
а $D$ --- нескалярной матрицей $H$ размера $2\times2$.
Можно одновременно потребовать
\[
 c_i=\gamma(C)_i\ne0\quad(i\in G):
\]
каждая координата на $\test_G$ ненулевая, а условие обратимости
и эти конечные условия задаются ненулевыми полиномами от двух
координат теста. Положим $t_i=\gamma(D)_i/c_i$.

Вычитая две формулы \eqref{eq:cycles} для $C,D$, получаем матричное
отношение с коэффициентами при $V_i$:
\begin{equation}\label{eq:cross-relation}
 \sum_{i\in G} c_i(t_iI-H)X_i\boxtimes V_i=0.
\end{equation}
Здесь равенство понимается покоординатно по матрице первого множителя;
оно также непосредственно следует из формулы
$\eta'(uC)=\sum_i\gamma(C)_i(uU_i\eta')V_i$ для любых $u,\eta'$.

По \eqref{eq:mass} $\dim\Span V_G\geqs6$.
Если $V_G$ независимы, каждый первый коэффициент в
\eqref{eq:cross-relation} нулевой. Тогда столбцы каждого $X_i$ лежат
в собственной прямой $H$ с собственным значением $t_i$.
У нескалярной матрицы $2\times2$ над замкнутым полем не больше двух
собственных направлений. На каждом по лемме~\ref{lem:growing-block}
не больше двух индексов, значит $|G|\leqs4$, противоречие.
Следовательно,
\[
 |G|=7,\qquad \dim\Span V_G=6.
\]
Все отношения между $V_i$ порождены одним ненулевым вектором $w$:
$\sum_iw_iV_i=0$. Каждая координата матричного отношения
\eqref{eq:cross-relation} кратна $w$. Поэтому существует матрица
$T$ формата $2\times3$ с
\begin{equation}\label{eq:resolvent}
 (t_iI-H)X_i=(w_i/c_i)T\qquad(i\in G).
\end{equation}

Докажем, что $\rk T=2$. Пусть $\lambda,\mu$ --- корни
характеристического многочлена $H$, возможно совпадающие. Оба сдвига
$H-\lambda I$, $H-\mu I$ ненулевые ранга один, и
\begin{equation}\label{eq:CH}
 (H-\lambda I)(H-\mu I)=0.
\end{equation}
Если $T=0$, предыдущий подсчёт снова даёт $|G|\leqs4$.
Пусть $T=ab$ ранга один. Если $a$ не собственный столбец $H$,
на параметрах $t_i=\lambda,\mu$ умножение \eqref{eq:resolvent} на
противоположный сдвиг даёт $w_i=0$ и помещает $X_i$ в соответствующую
собственную столбцовую прямую. Вне этих двух параметров сдвиг обратим,
и все $X_i$ имеют общее строковое направление $Kb$.
Каждый из трёх классов имеет мощность не больше двух; всего не больше
шести индексов, противоречие.

Если $Ha=\lambda a$, то умножение \eqref{eq:resolvent} на
$H-\lambda I$ и \eqref{eq:CH} показывают, что для $t_i\ne\mu$
все столбцы $X_i$ лежат в $Ka$. Но именно эти индексы образуют носитель
ненулевого теста $D-\mu C$. Их не больше двух по
лемме~\ref{lem:growing-block}, тогда как любой ненулевой тест требует
не меньше трёх индексов. Случай $Ha=\mu a$ аналогичен.
Таким образом, $\rk T=2$.

На корнях $t_i=\lambda,\mu$ левая часть \eqref{eq:resolvent} имеет
ранг не больше одного, поэтому $w_i=0$. Обозначим
$J_0=\{i\in G:t_i\ne\lambda,\ t_i\ne\mu\}$.
Для $i\in J_0$ из \eqref{eq:CH} и \eqref{eq:resolvent} следует
\[
 c_i(t_i-\lambda)X_i
 =w_iT+\frac{w_i}{t_i-\mu}(H-\lambda I)T.
\]
Умножим на билинейный второй множитель $V_i$ и просуммируем.
Первое слагаемое исчезнет, поскольку $\sum_iw_iV_i=0$ и $w_i=0$ вне $J_0$:
\[
 \sum_{i\in J_0}c_i(t_i-\lambda)(X_i\boxtimes V_i)
 =((H-\lambda I)T)\boxtimes
     \left(\sum_{i\in J_0}\frac{w_i}{t_i-\mu}V_i\right).
\]
Умножение первых матриц слева на $J$ переносит равенство к исходным
$U_i$. Выберем $j$ с $w_j\ne0$. Тогда $j\in J_0$, коэффициент при
его паре ненулевой и $\rk U_j\geqs\rk X_j\geqs\rk T=2$ по
\eqref{eq:resolvent}. Но первый множитель справа,
$J(H-\lambda I)T$, имеет ранг не больше одного.
Это противоречит лемме~\ref{lem:exchange}.
\end{proof}

\section{Четыре атома: остаток из восьми слагаемых}

\subsection{Пятимерный аннулятор}

\begin{lemma}\label{lem:annihilator}
Пусть $C_j=a_js_j$, $1\leqs j\leqs4$, --- матрицы ранга один $3\times3$,
причём любые три $a_j$ и любые три $s_j$ независимы.
Тогда пространство
\[
 L=\{T\in\Mat_3(K):\tr(TC_j)=0\ (1\leqs j\leqs4)\}
\]
имеет размерность пять и обладает базисом из матриц ранга один.
\end{lemma}
\begin{proof}
Независимые замены базисов двух внешних пространств и ненулевые
масштабирования позволяют выбрать столбцовые представители в виде
\[
 e_1,\quad e_2,\quad e_3,\quad e_1+e_2+e_3.
\]
Строковые представители независимо приводятся к строкам с теми же
координатами. Поэтому уравнения аннулятора имеют вид
\[
 T_{11}=T_{22}=T_{33}=0,\qquad \sum_{i\ne j}T_{ij}=0.
\]
Это четыре независимых уравнения. Рассмотрим шесть матриц ранга один:
\[
 E_{12}-E_{13},\ E_{21}-E_{23},\ E_{31}-E_{32},\quad
 E_{21}-E_{31},\ E_{12}-E_{32},\ E_{13}-E_{23}.
\]
Они порождают указанную гиперплоскость в шестимерном пространстве
внедиагональных координат: это разности соседних координат вдоль
связного шестичленного цикла. Любые пять, оставляющие связный путь,
образуют базис. Обратные замены базисов сохраняют ранги.
\end{proof}

\begin{prop}\label{prop:residual-model}
Если в алгоритме длины двадцать есть четыре атомарных теста с различными
носителями, то существуют пятимерное пространство $L$, его базис
$b_1,\ldots,b_5$, восемь векторов $w_i\in L$ и матрицы
$A_{ji}\in\Mat_3(K)$ ранга не больше одного, удовлетворяющие
\begin{equation}\label{eq:reconstruction}
 \sum_{i=1}^{8}\tr(XA_{ji})w_i=\tr(X)b_j
            \qquad\text{для всех }X\in\Mat_3(K),\quad 1\leqs j\leqs5.
\end{equation}
Систему равенств \eqref{eq:reconstruction} далее называем
\emph{реконструкцией выходов} $b_1,\ldots,b_5$.
\end{prop}
\begin{proof}
По лемме~\ref{lem:atom} четыре носителя непересекающиеся и содержат
двенадцать индексов. Пусть $J$ --- восемь оставшихся индексов,
$L$ --- аннулятор четырёх тестов. Последние удовлетворяют условиям
леммы~\ref{lem:annihilator}; выберем её базис $b_j$ ранга один.
Для $i\in J$ имеем $W_i\in L$; положим $w_i=W_i$,
$A_{ji}=V_i b_j U_i$. Эти матрицы имеют ранг не больше одного.
В \eqref{eq:dual-sandwich} с $T=b_j$ атомные слагаемые равны нулю:
если $U_i=ab_i$, $V_i=c_i s$, то
$V_i b_j U_i=c_i(s b_j a)b_i=0$, поскольку
$s b_j a=\tr(b_j as)=0$. Остаётся \eqref{eq:reconstruction}.
\end{proof}

\subsection{Координаты трёхмерного ядра}

\begin{lemma}\label{lem:kernel-model}
Если выполнено \eqref{eq:reconstruction}, то существуют восемь попарно
непропорциональных ненулевых столбцов $\mu_i\in K^3$, порождающих $K^3$,
и для каждого $j$ линейное отображение $Z_j:K^3\to\Mat_3(K)$, для которых
\[
 A_{ji}=a_{ji}I+Z_j(\mu_i),\qquad (Z_j(x))_{11}=0.
\]
Кроме того, $I\in\Span_i A_{ji}$ и
\[
 3\leqs\dim\Span_i A_{ji}=1+\rk Z_j\leqs4.
\]
\end{lemma}
\begin{proof}
Положим $F:K^8\to L$, $F(c)=\sum_i c_iw_i$. Подстановка $X=E_{11}$
в \eqref{eq:reconstruction} показывает, что все $b_j$ лежат в образе $F$.
Значит, $F$ сюръективно и $\dim\ker F=3$.
Возьмём $\mu_i$ --- координаты ограничений $i$-х координатных форм
на $\ker F$ в некотором базисе $(\ker F)^*$.
Они порождают всё это двойственное пространство.

Если $\mu_i,\mu_k$ зависимы, ненулевая форма
$a\varepsilon_i+b\varepsilon_k$ зануляет $\ker F$ и имеет вид
$\ell\circ F$ для $\ell\in L^*$. Применение $\ell$ к
\eqref{eq:reconstruction} даёт
\[
 aA_{ji}+bA_{jk}=\ell(b_j)I.
\]
Левая часть имеет ранг не больше двух, поэтому $\ell(b_j)=0$
для всех $j$. Тогда $\ell=0$ и $a=b=0$, противоречие.
Это также исключает нулевые $\mu_i$.

Положим $a_{ji}=(A_{ji})_{11}$. Для каждой клетки матрицы семейство
координат $A_{ji}-a_{ji}I$ лежит в $\ker F$, по
\eqref{eq:reconstruction}. Следовательно, эти матрицы получаются
линейным отображением $Z_j$ от координат $\mu_i$; их клетка $(1,1)$ нулевая.

Выберем $\ell$ с $\ell(b_j)=1$ и снова применим её к реконструкции:
$I=\sum_i\ell(w_i)A_{ji}$. Поэтому
$\Span_i A_{ji}=KI\oplus\im Z_j$.
Его размерность не меньше трёх: выбрав базис из матриц ранга не больше
одного, нельзя представить единицу ранга три менее чем тремя матрицами.
Верхняя граница четыре следует из $\dim K^3=3$.
\end{proof}

\subsection{Квадратичное отделение и конкурентные прямые}

\begin{definition}
Ненулевой столбец $\mu\in K^3$ задаёт проективную точку $[\mu]$;
пропорциональные столбцы задают одну точку. Проективная прямая ---
множество одномерных подпространств некоторой плоскости в $K^3$.
Три различные прямые называются \emph{конкурентными}, если проходят
через одну точку, называемую их центром.
В покрытии семейства точек такими прямыми \emph{активная группа}
прямой --- индексы точек на ней вне центра.

Индекс $i$ называется \emph{квадратично отделимым} в семействе
$[\mu_1],\ldots,[\mu_8]$, если существует однородный квадратичный
полином $q$ трёх переменных с $q(\mu_i)\ne0$ и $q(\mu_k)=0$ при $k\ne i$.
\end{definition}

\begin{lemma}\label{lem:separator}
Пусть матрицы $A_i$ ранга один попарно непропорциональны,
$I\in\Span A_i$ и $\dim\Span A_i\leqs4$.
Тогда существуют индекс $i$ и линейная форма $\chi$ на матрицах с
\[
 \chi(I)=1,\quad\chi(A_i)\ne0,\quad\chi(A_k)=0\ (k\ne i).
\]
Кроме того, существует однородная квадратичная форма $q$ на матрицах,
инвариантная при добавлении скалярной единицы, которая отделяет тот же индекс.
\end{lemma}
\begin{proof}
Разобьём семейство $A_i$ на компоненты связности; их оболочки образуют
прямую сумму. Разложим $I=\sum_H T_H$ по этой сумме и положим
$d_H=\dim\Span A_H$, $r_H=\rk T_H$.
Лемма~\ref{lem:connected} даёт $2r_H\leqs d_H+1$: столбцы $T_H$ лежат
в оболочке первых факторов матриц $A_i$, $i\in H$, а строки ---
в оболочке вторых, так что обе размерности не меньше $r_H$.
Если $d_H\geqs2$, целочисленность даёт $3r_H\leqs2d_H$.
Если бы при каждом $d_H=1$ было $T_H=0$, получили бы
\[
 3=\rk I\leqs\sum_H r_H\leqs\frac23\sum_H d_H\leqs\frac83,
\]
противоречие. Значит, есть одномерная компонента с $T_H\ne0$.
Из попарной непропорциональности она состоит из одного $A_i$.
Проекция на эту компоненту, продолженная до линейной формы на всех
матрицах и нормированная значением на $I$, даёт $\chi$.

Матрица $A_i-\chi(A_i)I$ имеет ранг не меньше двух:
иначе из $\rk A_i=1$ следовало бы $\rk I\leqs2$.
Выберем её ненулевой минор порядка два. Этот минор, взятый у
$T-\chi(T)I$, определяет квадратичную форму $q(T)$.
При $k\ne i$ она равна соответствующему минору $A_k$ и потому нулевая.
Также $q(T+cI)=q(T)$ вследствие $\chi(I)=1$.
\end{proof}

\begin{lemma}\label{lem:slice-dichotomy}
В модели леммы~\ref{lem:kernel-model} для каждого $j$ имеет место
хотя бы одна из двух возможностей:
\begin{enumerate}
\item существует квадратично отделимый индекс $i$ с $b_j\in Kw_i$;
\item восемь точек $[\mu_i]$ покрываются тремя конкурентными прямыми,
на каждой из которых есть хотя бы две точки вне центра;
при этом $b_j$ лежит в оболочке $w_i$ каждой из трёх групп индексов,
лежащих на соответствующей прямой вне центра.
\end{enumerate}
\end{lemma}
\begin{proof}
Фиксируем индекс $j$ и далее его опускаем.
Если $\dim\Span A_i=4$, отображение $Z$ инъективно.
Поэтому $A_i$ ненулевые и попарно непропорциональны: соответствующее
соотношение после вычитания клетки $(1,1)$ дало бы такое же
соотношение между $\mu_i$. Применим лемму~\ref{lem:separator}.
Из реконструкции и $\chi(I)=1$ получаем
\[
 b=\sum_k\chi(A_k)w_k=\chi(A_i)w_i.
\]
Квадратичная форма $x\mapsto q(Z(x))$ отделяет $\mu_i$, поскольку
\[
 q(Z(\mu_k))=q(A_k-a_kI)=q(A_k).
\]

Пусть $\dim\Span A_i=3$. Выберем в этом пространстве базис из матриц
ранга один. Их линейная комбинация равна $I$, так что оба набора
факторов являются базисами и, после переноса ненулевых коэффициентов,
двойственны. Поэтому сопряжением все $A_i$ одновременно
диагонализуются; каждая ненулевая $A_i$ имеет ровно одну ненулевую
диагональную координату.

После этого сопряжения рассмотрим
\[
 \pi(x)=\bigl((Z(x))_{22}-(Z(x))_{11},\
                 (Z(x))_{33}-(Z(x))_{11}\bigr).
\]
Её ядро равно $\ker Z$: равенство трёх диагоналей означает скалярную
матрицу, а исходное условие $(Z(x))_{11}=0$ исключает ненулевой скаляр.
Следовательно, $\rk\pi=2$, $\dim\ker\pi=1$.
Три координатные матричные прямые после факторизации по $KI$
дают три различные прямые в $K^2$. Их прообразы при $\pi$ ---
три плоскости через $\ker\pi$, то есть три конкурентные проективные
прямые. Они покрывают все $[\mu_i]$.

Обозначим $a$-ю диагональную координату сопряжённой $A_i$ через $c_{ai}$.
Применяя диагональную форму к реконструкции, получаем
$b=\sum_i c_{ai}w_i$ для каждого $a=1,2,3$.
Ненулевые $c_{ai}$ лежат на соответствующей прямой вне центра.
Каждая активная группа непуста. Если во всех трёх группах хотя бы по
два индекса, получаем второй случай. Иначе группа состоит из одного
индекса $i$, и $b\in Kw_i$. Произведение линейных уравнений двух
остальных прямых --- квадратичный полином, ненулевой на $\mu_i$
и нулевой на остальных точках. Получен первый случай.
\end{proof}

\begin{lemma}\label{lem:projective-count}
Для восьми различных проективных точек верно следующее.
\begin{enumerate}
\item Квадратично отделимых индексов не больше трёх.
\item Если есть покрытие тремя конкурентными прямыми, на каждой из
которых хотя бы две точки вне центра, то квадратично отделимых индексов нет.
\item Различных центров таких покрытий не больше четырёх.
\end{enumerate}
\end{lemma}
\begin{proof}
\emph{1.} Значения квадратичных форм образуют подпространство $K^8$
размерности не больше шести. Если отделимых индексов $s$, их
координатные векторы принадлежат этому пространству. На оставшихся
точках ранг вычисления квадратичных форм не меньше
$\min(3,8-s)$: любые три различные проективные точки отделяются
произведениями линейных форм. Значит,
$s+\min(3,8-s)\leqs6$, что влечёт $s\leqs3$.

\emph{2.} Пусть квадратичная форма зануляется на всех точках, кроме,
возможно, одной. Её ограничение на проективную прямую, если оно
ненулевое, имеет не больше двух нулей. Поэтому три нуля на прямой
вынуждают эту прямую быть линейным множителем формы.
Если центр не входит в восемь точек, после удаления одной точки
остаются семь точек на трёх прямых, с непустыми группами, из которых
как минимум две имеют размер не меньше двух. Можно выбрать на них
соответственно три, две и одну точку. Первая тройка вынуждает один
линейный множитель; вторая пара --- второй. Последняя точка, лежащая
вне этих двух прямых, даёт противоречие.
Если центр входит в семейство и удаляется именно он, остаются семь
внешних точек, не менее двух на каждой прямой, и тот же довод применим.
Если удалена внешняя точка, центр остаётся нулём. На двух прямых
осталось хотя бы по две внешние точки; вместе с центром это две
тройки. Форма должна быть произведением уравнений этих двух прямых,
но на третьей осталась внешняя точка, где оно ненулевое.

\emph{3.} Если некоторая прямая содержит хотя бы четыре точки,
она обязана входить в каждое тройное покрытие: иначе три прямые
покрытия пересекали бы её не более чем в трёх точках.
Все центры лежат на ней. При двух различных центрах вне неё может
быть не более четырёх точек, по четырём пересечениям оставшихся пар
прямых; активность требует не меньше четырёх.
Зафиксируем внешнюю точку. В каждом покрытии на её прямой есть
другая внешняя точка; возможны лишь три партнёра. Партнёр определяет
прямую и её пересечение с исходной прямой, то есть центр. Поэтому
центров не больше трёх.

Остаётся случай, когда на любой прямой не больше трёх точек.
Каждое покрытие имеет хотя бы две прямые с тремя точками.
Два различных центра не могут использовать одну такую прямую:
иначе пять внешних точек лежали бы в четырёх пересечениях остальных
пар прямых. Если бы центров было пять, получили бы десять различных
трёхточечных прямых. Каждая содержит три неупорядоченные пары точек,
и никакая пара не повторяется. Но $10\cdot3>\binom82=28$.
\end{proof}

\begin{lemma}\label{lem:same-center}
В модели \eqref{eq:reconstruction} два независимых выхода $b,b'$
не могут иметь покрытия второго типа леммы~\ref{lem:slice-dichotomy}
с одним центром.
\end{lemma}
\begin{proof}
Для фиксированного центра каждая внешняя точка однозначно определяет
проходящую через неё прямую. Поэтому активные группы обоих покрытий
совпадают с точностью до перестановки.
Пусть $E_1,E_2,E_3\subseteq K^8$ --- координатные пространства этих
трёх непересекающихся групп. Если $H=F(E_1)\cap F(E_2)\cap F(E_3)$,
то $2\dim H\leqs\dim\ker F=3$.
Действительно, выберем линейные $s_a:H\to E_a$ с $F(s_a(h))=h$;
отображение
$(h,k)\mapsto s_1(h)-s_2(h)+s_1(k)-s_3(k)$ инъективно
и принимает значения в $\ker F$. Инъективность видна по второй и
третьей непересекающимся группам координат.
Следовательно, $\dim H\leqs1$, и два независимых выхода в $H$ невозможны.
\end{proof}

\begin{theorem}\label{thm:no-eight}
Реконструкция \eqref{eq:reconstruction} пяти независимых выходов
восемью слагаемыми ранга не больше одного невозможна над любым полем.
\end{theorem}
\begin{proof}
Построим общее семейство восьми точек по лемме~\ref{lem:kernel-model}.
Если все пять выходов имеют первый тип из
леммы~\ref{lem:slice-dichotomy}, они лежат в оболочке не более трёх
векторов $w_i$ по пункту 1 леммы~\ref{lem:projective-count},
что противоречит независимости.
Значит, есть хотя бы одно покрытие второго типа. По пункту 2
квадратично отделимых индексов теперь нет, так что все пять выходов
имеют второй тип. По лемме~\ref{lem:same-center} их центры различны,
вопреки пункту 3 леммы~\ref{lem:projective-count}.
\end{proof}

\begin{cor}\label{cor:no-four}
Алгоритм длины двадцать не может иметь четыре атомарных теста
с различными носителями.
\end{cor}
\begin{proof}
Предложение~\ref{prop:residual-model} сводит такой алгоритм к
реконструкции, запрещённой теоремой~\ref{thm:no-eight}.
\end{proof}

\section{Три атома и завершение доказательства}

\begin{prop}\label{prop:three}
Для нормированного алгоритма длины двадцать равенства $MD=D$ и
$\dim D=3$ невозможны.
\end{prop}
\begin{proof}
По лемме~\ref{lem:component-cost} пространство $D$ порождается общими
собственными векторами умножений на $M$. Выберем три независимых
таких вектора. Лемма~\ref{lem:stationary} превращает их в атомарные
тесты $C_j=a_js_j$ с различными носителями.
По лемме~\ref{lem:atom} столбцы $a_j$ образуют базис, как и строки $s_j$.
Пусть $(u_j)$ --- базис строк, двойственный $(a_j)$.

Поскольку $\gamma(C_j)$ --- общий собственный вектор,
$C_{\gamma(C_j)}(u)$ пропорциональна $(ua_j)C_j$ с ненулевым
постоянным множителем. Закон \eqref{eq:compatibility} даёт
\[
 (ua_j)s_j B_z(v)=(va_j)s_j B_z(u).
\]
Выбирая $u=u_k$, получаем $s_jB_z(u_k)=0$ при $j\ne k$.
Следовательно,
\[
 B_z(u_k)=s_k^*f_k(z)
\]
для некоторых линейных $f_k:Z\to F$. Каждое $f_k$ инъективно:
его нуль означал бы $B_z(u_k)=0$, а $u_k\ne0$.
По лемме~\ref{lem:injective-product} $r\geqs9z+3$.
Но \eqref{eq:dz} при $d=3$, $r=20$ даёт $z\geqs2$,
откуда $r\geqs21$, противоречие.
\end{proof}

\begin{proof}[Доказательство основной теоремы]
Предположим, что над исходным полем существует алгоритм длины
$r\leqs20$. Расширим поле до алгебраического замыкания и удалим
нулевые слагаемые. Предложение~\ref{prop:critical} показывает, что
оставшаяся длина равна двадцати.

Выберем среди алгоритмов этой длины минимальный по первой сумме
рангов, затем нормируем его по лемме~\ref{lem:generic}.
Предложение~\ref{prop:no-growth} даёт $MD=D$, а
предложение~\ref{prop:critical} --- $d\geqs3$.
Если $d\geqs4$, выберем четыре независимых общих собственных вектора
в $D$; по леммам~\ref{lem:component-cost} и~\ref{lem:stationary}
это четыре атомарных теста. Их носители различны, поскольку атомы
с одинаковым носителем пропорциональны. Получено противоречие со
следствием~\ref{cor:no-four}.
Поэтому $d=3$, но этот случай исключён предложением~\ref{prop:three}.

Алгоритма длины не больше двадцати нет и над исходным полем.
Следовательно, $R_K\langle3,3,3\rangle\geqs21$.
\end{proof}
\end{document}
